\captionsetup[table]{
    justification=raggedleft,
    singlelinecheck=false,
    font=it,
    labelfont=it,
    skip=10pt
}

\newcommand{\uniterablesection}[1]{%
    \clearpage\phantomsection
    \addcontentsline{toc}{section}{\MakeUppercase{#1}}
    \section*{#1}
}

% Титульна сторінка
\newcommand{\maketitlepage}[9]{%
    \begin{titlepage}
    \thispagestyle{empty}
    \begin{center}
    \MakeUppercase{\textbf{#1}}\\[0.5em] \textbf{#2}\\[1em] \textbf{#3}\\[8em]
    \MakeUppercase{\textbf{#4}}\\[0em] з дисципліни «#5»\\[0em] на тему:\\[0em] «#6»
    \end{center}
    \vfill\noindent\hfill
    \begin{tabular}{l} Виконав студент ННІКІТІ: \\ групи: #7 \\ #8 \\[1em] Науковий керівник: \\ \parbox[t]{8cm}{#9} \end{tabular}
    \vspace{6em}\begin{center} Рівне – 2026 \end{center}
    \end{titlepage}
}

% Оформлення рисунків
\DeclareCaptionLabelFormat{myformat}{\textbf{#1 #2.}}

\captionsetup[table]{
    labelsep=period,
    justification=raggedleft,
    singlelinecheck=false,
    font=it
}

\captionsetup[figure]{
    name=Рис.,
    labelformat=myformat,
    labelsep=space,
    justification=centering,
    singlelinecheck=false,
    font=normalsize,
    textfont=normalfont,
}

\counterwithin{figure}{section}
\counterwithin{table}{section}

\renewcommand{\thefigure}{\arabic{section}.\arabic{figure}}
\renewcommand{\thetable}{\arabic{section}.\arabic{table}}
\renewcommand{\thesection}{\Roman{section}}
\renewcommand{\thesubsection}{\arabic{section}.\arabic{subsection}}
\renewcommand{\thesubsubsection}{\arabic{section}.\arabic{subsection}.\arabic{subsubsection}}

\newcommand{\setfigure}[4][]{
    \begin{center}
    \nopagebreak
    \includegraphics[width=#2]{#3}
    \captionof{figure}{#4}
    \label{#1}
    \end{center}
    \vspace{-1em}
}

% --- customSection ---
\newcommand{\customSection}[3][]{%
    \ifx
        \relax#3\relax\else
        \vspace{#3}
    \fi
    \section{#2}\label{#1}%
}

% --- customSubsection ---
\newcommand{\customSubSection}[3][]{%
    \ifx
        \relax#3\relax\else
        \vspace{#3}
    \fi
    \subsection{#2}\label{#1}%
}

% --- customSubSubSection ---
\newcommand{\customSubSubSection}[3][]{%
    \ifx
        \relax#3\relax\else
        \vspace{#3}
    \fi
    \subsubsection{#2}\label{#1}%
}

\newcommand{\customBibliography}[1]{%
    \phantomsection
    \label{#1}
    \begin{enumerate}[noitemsep, topsep=0pt, leftmargin=*, label=\arabic*.]
}

\newcommand{\closeCustomBibliography}{%
    \end{enumerate}
}

% Книга або стаття (IEEE Style)
% #1 - Автори (І. Прізвище)
% #2 - Назва (Книги або Статті)
% #3 - Видавництво / Конференція / Місто
% #4 - Рік, сторінки, DOI
\newcommand{\bibitemieee}[4]{%
    \item #1, ``#2,'' \textit{#3}, #4.
}

% Електронний ресурс (IEEE Style)
\newcommand{\bibitemweb}[3]{%
    \item #1. [Online]. Available: \url{#2}. [Accessed: #3].
}

% коротке посилання в тексті
\newcommand{\citecustom}[1]{[#1]}

% посилання на таблицю
\newcommand{\tabref}[1]{табл.~\ref{#1}}
